\par\medskip\noindent\begin{minipage}{\linewidth}\centering
\begin{tikzpicture}\begin{axis}[width=.88\linewidth,height=6.4cm,axis lines=middle,grid=major,grid style={gray!25},tick label style={font=\normalsize},label style={font=\normalsize},clip=false,xmin=-.2,xmax=3.4,ymin=-1,ymax=10,xlabel={$x$},ylabel={$y$},xtick={0,1,2,3},ytick={0,2,4,6,8},legend style={at={(.02,.98)},anchor=north west,fill=white,font=
ormalsize}]\addplot[bleu,very thick,domain=0:3.1]{x^2};\addlegendentry{$y=x^2$}\addplot[vert,very thick,dashed,domain=.8:3.2]{4*x-4};\addlegendentry{$y=4x-4$}\addplot[only marks,mark=*] coordinates {(1,0)(2,4)};\draw[dotted] (axis cs:1,0)--(axis cs:2,0)--(axis cs:2,4);\node[above left,fill=white,inner sep=2pt]at(axis cs:1,0){$(1;0)$};\node[right,fill=white,inner sep=2pt]at(axis cs:2,4){$A(2;4)$};\end{axis}\end{tikzpicture}
\par\medskip La tangente en A passe aussi par $(1;0)$. Sa pente vaut $\dfrac{4-0}{2-1}=4$. Le tracé discontinu permet de la distinguer de la courbe, même en noir et blanc.\end{minipage}\par\medskip
